\section{Подготовка реакции на перегрузку}
Метод объединяет исполнимую модель процесса с обученными зависимостями поведения, сравнивает допустимые вмешательства при заданных нагрузках и проверяет сохранённое правило на независимых сценариях. Результат — действие с наблюдаемым условием активации и зафиксированными предпосылками выбора.

\subsection{Исполнение процесса и обратные связи}
Экземпляр проходит операции, описанные моделью $\mathcal M=(N,\mathcal R,K,\alpha,F_D,F_R)$: структура $N$, ресурсные пулы $\mathcal R$, мощности $K$, назначение операций пулам $\alpha$, функции длительности и маршрутизации. Поддерживаются блочные последовательности, XOR-развилки, AND-разветвления с синхронизацией и вложенные циклы с наблюдаемыми операциями. При постановке операции в очередь резервируются входные токены; выходные появляются после завершения. Готовые операции используют конечные FIFO-пулы с занятостью $0\leq b_r(t)\leq K_r$. Незавершённая работа $w(t)$ учитывает каждый экземпляр один раз, в том числе при нескольких одновременных операциях.

Журнал жизненного цикла содержит поступление экземпляра, моменты готовности, начала и завершения операций, их тип, ресурсный пул и наблюдаемые атрибуты экземпляра. Длительность выбирается при начале обслуживания $t_s$:
\begin{equation}\label{eq:duration}
\log D_a=f_a(x,\ell(t_s))+\epsilon_a,\qquad \epsilon_a\sim\mathcal N(0,\sigma_a^2).
\end{equation}
Для бинарного решения о возврате в момент завершения $t_c$:
\begin{equation}\label{eq:return}
R\sim\mathrm{Bernoulli}\!\left(\operatorname{sigmoid}(g(x,\ell(t_c)))\right).
\end{equation}
Здесь $x$ содержит наблюдаемый класс и, при конкуренции процессов, идентификатор процесса. Нагрузочные признаки включают $\log(1+q_a/K_{\alpha(a)})$, $\log(1+w/S)$ и WIP отдельного процесса. Масштаб $S$ фиксируется до оценки. Очередь при старте включает начинающуюся операцию; признаки в момент завершения не включают предполагаемый возврат. Будущие состояния не используются как входные признаки.

При нескольких маршрутах отдельная мультиномиальная модель для каждой исходной операции предсказывает следующую готовую операцию либо завершение экземпляра. Для одновременно готовых операций целевая метка определяется фиксированным лексическим порядком. Вес исполнимой альтернативы равен сумме предсказанных вероятностей достижимых следующих меток; вероятности бинарного цикла затем нормируются. Класс экземпляра различает неоднородные маршруты, а отдельные исходные операции — разные точки возврата. Для длительности используется ridge-регрессия с параметром $\alpha=1$, для маршрутизации — логистическая регрессия с $C=10$. Нагрузочные признаки маршрутизации добавляются в заданных точках вмешательства в повторную обработку.

Статические модели сохраняют атрибуты и ресурсные ограничения, но исключают нагрузочные признаки. Retained-load сохраняет зависимости обслуживания и маршрутизации. Предиктивный селектор оставляет компонент, только если улучшение вероятностной функции потерь на валидации превышает $0.01$ для длительности или $0.005$ для маршрутизации. Используются отрицательное логарифмическое правдоподобие нормального распределения логарифма длительности и кросс-энтропия маршрутов. Сравнение проверяет, сохраняет ли предиктивный отбор информацию, необходимую для последующего выбора реакции. Исходная последовательная реализация P0 сохраняет прежние масштабы признаков, коэффициенты и результаты; расширение на сложные структуры использует отдельные модели операций и точек возврата.

Структура задаётся заранее либо восстанавливается по обучающим трассам. Обнаружение сравнивает стандартный Inductive Miner с кандидатом, который проверяет разложение строгого цикла перед более свободным выделением параллельности. Второй вариант также выводит бинарные условия по классу или процессу, известным при поступлении и постоянным внутри экземпляра, если каждому значению атрибута соответствуют не менее 10 обучающих экземпляров с постоянным и различающимся набором операций. Среди исполнимых кандидатов, допускающих все обучающие трассы по оптимальным выравниваниям, выбирается наибольшая обучающая precision, затем наименьшее число переходов. Ресурсные отображения, потенциальные точки возврата, соответствующие им наблюдаемые очереди и семантика действий задаются инженером; истинная структура переходов и тестовые исходы не участвуют в выборе.

Исполнение должно сохранять смысл точки вмешательства. Модель следующей операции не различает альтернативы, начинающиеся одинаковым наблюдаемым продолжением. Например, $\mathrm{seq}(\mathrm{loop}(B,D),D)$ и $\mathrm{loop}(\mathrm{seq}(B,D),\tau)$ допускают одинаковые наблюдаемые трассы, но во втором случае решение о повторе принимается после D, где можно представить его исход и эффект вмешательства. Компилятор выполняет это преобразование и отклоняет оставшиеся развилки с пересекающимися наборами следующих меток. Условия по наблюдаемым атрибутам соблюдаются при имитации обученной модели. Поэтому поддерживаемый подкласс требует детерминированных условий по наблюдаемым атрибутам либо различимых следующих меток после преобразования. Истинный генератор не используется для исправления обнаруженной структуры.

\subsection{Сравнение допустимых реакций при заданных нагрузках}
Инженер задаёт библиотеку сценариев $\mathcal S_p$, действия $\mathcal A$, стоимости и кандидатные пороги WIP. Реакция $\pi=(a,q)$ активируется после двух последовательных превышений:
\begin{equation}\label{eq:trigger}
\tau_q=\inf\{t_k:w(t_{k-1})\geq q\ \land\ w(t_k)\geq q\}.
\end{equation}
Правило использует наблюдаемый WIP без знания будущих поступлений. Дополнительный работник влияет на будущие старты, а сокращение rework — на будущие возвраты. Уже выбранные длительности обслуживания не меняются. При горизонте $H$ целевая функция равна
\begin{equation}\label{eq:cost}
J(\pi)=\int_0^H w_\pi(t)\,dt+c_a(H-\tau_q)_+,
\end{equation}
причём стоимость действия равна нулю, если активация не произошла. Интеграл WIP вычисляется методом трапеций по наблюдениям. Поддержание дополнительного работника стоит $2$ за единицу времени, сокращение rework — $0.5$. Это эквиваленты времени незавершённых экземпляров, а не денежные оценки. Терминальный остаток работы сообщается отдельно.

Планировщик выбирает минимальную среднюю имитационную стоимость по заданным сценариям и повторениям:
\begin{equation}\label{eq:policy}
\widehat\pi=\arg\min_{\pi\in\Pi}\frac{1}{|\mathcal S_p|L}
\sum_{s\in\mathcal S_p}\sum_{l=1}^{L}J_{\widehat{\mathcal M},s,l}(\pi).
\end{equation}
Все варианты моделей используют одинаковые действия, пороги, сценарии и стоимости. Равенства разрешаются фиксированным порядком каталога. Случайные потоки, индексированные счётчиками, обеспечивают парность сравнений действий. Выбранная политика сохраняется до генерации независимых тестовых поступлений.

\subsection{Проверка сохранённой реакции}
Тест измеряет реализованную стоимость, частоту длительной перегрузки, активацию, время предупреждения и длительности завершённых экземпляров. Время предупреждения определяется относительно начала тяжёлой перегрузки в парной траектории без вмешательства; этот контрфактический исход используется только при оценке. Предотвращение события отличается от снижения его интегральной стоимости. Реакция может уменьшить очередь, сохранив выполнение критерия тяжёлого события.

Запись о реакции связывает действие и порог с обученной моделью, планировочными нагрузками, наблюдавшимся диапазоном, предположениями о вмешательствах и тестовыми исходами. Эти сведения нужны для эксплуатационного согласования и определения необходимости новых измерений или калибровки.
